\section*{Capítulo \ref{ch:b1:elimination} --- \texorpdfstring{$\beta$}{Beta}-eliminação}

\begin{solution}{exo:b1:elimination:1}
2-Bromobutano: but-1-eno, \cip{E}- e \cip{Z}-but-2-eno; o majoritário é o
\cip{E}-but-2-eno. 2-Bromo-2-metilbutano: 2-metilbut-2-eno (majoritário,
trissubstituído) e 2-metilbut-1-eno. Bromociclo-hexano: só ciclo-hexeno.
\end{solution}

\begin{solution}{exo:b1:elimination:2}
E1: $v = k[\mathrm{RBr}]$; E2: $v = k[\mathrm{RBr}][\ce{EtO-}]$. Meça a
\omterm{def:b1:rate-laws:initial-rate}{velocidade inicial} de formação do alceno em várias concentrações de
etóxido: ela cresce proporcionalmente para a E2 e não muda para a E1.
\end{solution}

\begin{solution}{exo:b1:elimination:3}
Ao longo de C2--C3, o Br de C2 pode ficar em anti com qualquer um dos
hidrogênios de C3 ou com o grupo metila C4: as duas \omterm{def:b1:stereochemistry-in-depth:isomers}{conformações} com um H
em anti ao Br podem reagir por E2 (dando os alcenos \cip{E} e \cip{Z});
a terceira não pode eliminar em direção a C3.
\end{solution}

\begin{solution}{exo:b1:elimination:4}
SN2 (primário, hidróxido, temperatura baixa); E2 (terciário, \omterm{def:b1:predominant-reaction:strong-acid}{base forte}
e volumosa); SN2 (iodeto, bom \omterm{def:b1:electronic-effects:nucleophile}{nucleófilo} e \omterm{def:b1:predominant-reaction:strong-acid}{base fraca}, \omterm{def:b1:intermolecular-forces-solvents:solvent-classes}{solvente
aprótico}); SN1 e E1, com a E1 aumentando com o aquecimento.
\end{solution}

\begin{solution}{exo:b1:elimination:5}
C3 leva um único hidrogênio. Em cada \omterm{def:b1:stereochemistry-in-depth:enantiomers}{diastereoisômero}, pôr esse H em anti
ao Br de C2 fixa as posições dos outros grupos: o metila de C2 e o etila
de C3 ficam do mesmo lado em um \omterm{def:b1:stereochemistry-in-depth:enantiomers}{diastereoisômero} e em lados opostos no
outro. Um dá o \cip{E}-, o outro o \cip{Z}-3-metilpent-2-eno: uma \omterm{def:b1:nucleophilic-substitution:stereospecific}{reação
estereoespecífica}.
\end{solution}

\begin{solution}{exo:b1:elimination:6}
O OH é protonado, \ce{R-OH2+}; a água sai, dando o \omterm{def:b1:electronic-effects:intermediates}{carbocátion} terciário
\ce{(CH3)2C+-CH2CH3}; uma base (água, \ce{HSO4-}) retira um próton de um
carbono vizinho. Produto majoritário (Zaitsev): 2-metilbut-2-eno.
\end{solution}

\begin{solution}{exo:b1:elimination:7}
O grupo terc-butila trava a \omterm{def:b1:stereochemistry-in-depth:conformations}{cadeira} em que ele fica \omterm{def:b1:stereochemistry-in-depth:conformations}{equatorial}. No
isômero \textit{trans}, o bromo também fica então \omterm{def:b1:stereochemistry-in-depth:conformations}{equatorial}, nunca em
anti a uma C--H: a E2 precisa esperar a \omterm{def:b1:stereochemistry-in-depth:conformations}{cadeira} invertida, muito
desfavorável. No isômero \textit{cis}, o bromo é \omterm{def:b1:stereochemistry-in-depth:conformations}{axial}, com dois
hidrogênios trans-diaxiais: E2 rápida.
\end{solution}

\begin{solution}{exo:b1:elimination:8}
O etóxido, pequeno, retira o hidrogênio que dá o alceno mais
substituído (Zaitsev). O terc-butóxido, volumoso, alcança mais
facilmente os hidrogênios expostos do grupo metila do que o hidrogênio
do \ce{CH2} interno: o alceno menos substituído.
\end{solution}

\begin{solution}{exo:b1:elimination:9}
O bromometano não tem carbono $\beta$. No 1-bromo-2,2-dimetilpropano, o
carbono $\beta$ é quaternário e não leva nenhum hidrogênio.
\end{solution}

\begin{solution}{exo:b1:elimination:10}
Os dois produtos precisam do \omterm{def:b1:electronic-effects:intermediates}{carbocátion}, formado na etapa lenta à
velocidade $k_1[\mathrm{RBr}]$. Ele se reparte então entre a captura
pelo etanol (\omterm{def:b1:rate-laws:order}{constante de velocidade} $k_S$) e a perda de um próton
(\omterm{def:b1:rate-laws:order}{constante de velocidade} $k_E$), ambas de pseudoprimeira ordem em
relação ao solvente: a fração de alceno é $k_E/(k_E + k_S)$,
independente da concentração do \omterm{def:b1:nucleophilic-substitution:substitution}{substrato}.
\end{solution}

\begin{solution}{exo:b1:elimination:11}
Partindo da \omterm{def:b1:stereochemistry-in-depth:isomers}{conformação} do \omterm{def:b1:stereochemistry-in-depth:enantiomers}{composto meso} com os dois Br em anti (os dois
fenilas ficam então em anti, os dois H em anti), gire C1 de $120^\circ$
para pôr seu H em anti ao Br de C2. Os dois grupos fenila ficam então do
mesmo lado do eixo H--Br: terminam em \textit{cis}. Em C1, Br tem
prioridade sobre Ph; em C2, Ph tem prioridade sobre H; o Br e o fenila
de C2 estão em \textit{trans}: \cip{E}-1-bromo-1,2-difenileteno.
\end{solution}

\begin{solution}{exo:b1:elimination:12}
Só a \omterm{def:b1:stereochemistry-in-depth:conformations}{cadeira} com o \omterm{def:b1:electronic-effects:nucleophile}{grupo de saída} \omterm{def:b1:stereochemistry-in-depth:conformations}{axial} reage: $v = k\,x\,[\mathrm{RY}][\mathrm{B}]$,
com $x$ a fração dessa \omterm{def:b1:stereochemistry-in-depth:conformations}{cadeira}. Cada grupo alquila \omterm{def:b1:stereochemistry-in-depth:conformations}{axial} que ela exige
custa vários kJ/mol (\qty{5.7}{kJ/mol} para um metila), dividindo $x$
por cerca de dez por grupo à temperatura ambiente: com dois ou três
grupos assim, $x$ cai para $10^{-2}$ ou $10^{-3}$ e a reação fica
igualmente mais lenta.
\end{solution}

\begin{solution}{pb:b1:elimination:1}
\textbf{1.} Uma \omterm{def:b1:stereochemistry-in-depth:conformations}{cadeira} com Cl, isopropila e metila todos \omterm{def:b1:stereochemistry-in-depth:conformations}{equatoriais}.
\textbf{2.} Isopropila e metila \omterm{def:b1:stereochemistry-in-depth:conformations}{equatoriais}, Cl \omterm{def:b1:stereochemistry-in-depth:conformations}{axial}.
\textbf{3.} \omterm{def:b1:stereochemistry-in-depth:enantiomers}{Diastereoisômeros}: eles diferem apenas em C1.
\textbf{4.} \omterm{def:b1:electronic-effects:nucleophile}{Grupo de saída} e hidrogênio $\beta$ trans-diaxiais.
\textbf{5.} Só na \omterm{def:b1:stereochemistry-in-depth:conformations}{cadeira} invertida, com Cl, isopropila e metila todos
\omterm{def:b1:stereochemistry-in-depth:conformations}{axiais}.
\textbf{6.} Os três grupos \omterm{def:b1:stereochemistry-in-depth:conformations}{axiais}: só o metila custa cerca de
\qty{5.7}{kJ/mol}, o grupo isopropila, maior, mais, e o cloro um pouco:
juntos, uns \qty{13}{kJ/mol}, logo uma fração de cerca de
$\exp(-13000/2479) = \num{5e-3}$, o valor que o problema adota.
\textbf{7.} C2 (um H, ele leva o grupo isopropila) e C6 (dois H).
\textbf{8.} Com o Cl \omterm{def:b1:stereochemistry-in-depth:conformations}{axial} em C1, o H \omterm{def:b1:stereochemistry-in-depth:conformations}{axial} de C2 (estando o isopropila
\omterm{def:b1:stereochemistry-in-depth:conformations}{equatorial}) e o H \omterm{def:b1:stereochemistry-in-depth:conformations}{axial} de C6: dois.
\textbf{9.} Na \omterm{def:b1:stereochemistry-in-depth:conformations}{cadeira} triaxial, o grupo isopropila ocupa a posição \omterm{def:b1:stereochemistry-in-depth:conformations}{axial}
de C2, de modo que seu H é \omterm{def:b1:stereochemistry-in-depth:conformations}{equatorial}: não está em anti. C6 tem um H
\omterm{def:b1:stereochemistry-in-depth:conformations}{axial}: em anti.
\textbf{10.} O cloreto de neomentila, dois; o cloreto de mentila, um.
\textbf{11.} Ao longo de C1--C6: Cl (\omterm{def:b1:stereochemistry-in-depth:conformations}{axial}, para cima) em C1 oposto ao H
\omterm{def:b1:stereochemistry-in-depth:conformations}{axial} (para baixo) em C6; as ligações do anel e o H \omterm{def:b1:stereochemistry-in-depth:conformations}{equatorial}
alternados entre eles.
\textbf{12.} Em uma \omterm{def:b1:stereochemistry-in-depth:conformations}{cadeira}, nenhuma ligação C--H é \omterm{def:b1:elimination:periplanar}{sinperiplanar} à
ligação C--Cl; forçá-la significaria um anel eclipsado e tensionado.
\textbf{13.} Retirando o H de C2: 4-metil-1-(propan-2-il)ciclo-hex-1-eno
(trissubstituído); retirando o H de C6: 3-metil-6-(propan-2-il)ciclo-hex-1-eno
(dissubstituído).
\textbf{14.} O alceno trissubstituído, o mais substituído e mais estável.
\textbf{15.} Os dois alcenos, \qty{78}{\percent} de trissubstituído:
coerente.
\textbf{16.} Só o alceno dissubstituído: o contrário da previsão de
Zaitsev, imposto pelo único hidrogênio anti disponível.
\textbf{17.} As duas são estereoespecíficas (exigência anti); o cloreto
de neomentila reage de modo regiosseletivo (78/22); o cloreto de mentila
dá um único \omterm{def:b1:stereochemistry-in-depth:isomers}{isômero constitucional}.
\textbf{18.} O carbono que leva o Cl é secundário e vizinho de um grupo
isopropila: congestionado para a SN2, enquanto o etóxido é uma \omterm{def:b1:predominant-reaction:strong-acid}{base
forte}.
\textbf{19.} $v \propto x\,n_{\mathrm{H}}$: fração de \omterm{def:b1:stereochemistry-in-depth:conformations}{cadeira} reativa
vezes o número de hidrogênios anti.
\textbf{20.} $0.95 \times 2 = 1.90$.
\textbf{21.} $\num{5.0e-3} \times 1 = \num{5.0e-3}$.
\textbf{22.} \textbf{$k(\text{neomentila})/k(\text{mentila}) = 1.90/\num{5.0e-3}
= 380$}.
\end{solution}
